% ===================================================================
%  ТАБЛИЦЯ № 13 — Готель. — Ресторан. — Кав'ярня.
%  Traduction 2026 d'apres texto/io, controlee sur texto/fr, texto/en
%  et texto/es. Consigne : voir texto/uk/10-tablycia-01.tex.
%
%  Deux notes, marquees toutes deux « (*) », et deux appels : celui du
%  bloc 01-2 (la chambre) et celui du bloc 09-1 (le menu). L'ordre les
%  departage.
%
%  DEUX CONDUCTEURS SUR LE MEME TROTTOIR, et il faut deux mots : celui
%  de l'omnibus (30) et celui du fiacre (92). L'ukrainien en a
%  justement deux, « кучер » et « візник », et le releve des objets les
%  garde distincts.
% ===================================================================

%%K t13-noto-1 noto
\VUnotes{143.2mm}{%
\textit{\VUgras{(*) Подробиці спальні див. у таблиці № 6.}}%
}

%%K t13-apar-1 apar
\VUsaut{3.0mm}
\VUtitre{15.0pt}{60}{ТАБЛИЦЯ № 13}
\VUsaut{1.6mm}
\VUpk{10.2pt}{40}{Готель. --- Ресторан.}
\VUsaut{2.0mm}
\VUpk{10.2pt}{40}{Кав'ярня.}
\VUsaut{3.4mm}

%%K t13-01-1 p
1. --- Приїхавши вчора ввечері до Руаяна, я спинився в \textit{Готелі Океану},
який мені порадив один приятель.

%%K t13-01-2 p
Мені дали прекрасну \VUgras{кімнату} \textsuperscript{(2)} на другому
\VUgras{поверсі} \textsuperscript{(1)}, де я чудово виспався (*).

%%K t13-02-1 p
2. --- Сьогодні вранці, вийшовши зі своєї кімнати, куди \VUgras{покоївка}
\textsuperscript{(3)} принесла мені сніданок, я зайшов до \VUgras{курильні}
\textsuperscript{(5)}, яка править і за \VUgras{читальню} \textsuperscript{(4)},
щоб проглянути \VUgras{газети} \textsuperscript{(6)}, покурюючи \VUgras{цигарку}
\textsuperscript{(8)}. Скінчивши читати, я спустився \VUgras{ліфтом}
\textsuperscript{(9)} і пішов прогулятися містом.

%%K t13-03-1 p
3. --- Що я забув спитати в готельного \VUgras{конторника}
\textsuperscript{(12)}, о котрій годині подають за \VUgras{спільним столом}
\textsuperscript{(11)}, я вернувся рано і скористався з нагоди взяти купіль у
готельній \VUgras{ванній} \textsuperscript{(10)}. Потім я пішов до \VUgras{каси}
\textsuperscript{(15)}, щоб подзвонити одному приятелеві. Але \VUgras{телефон}
\textsuperscript{(13)} був зайнятий; бачачи мій клопіт, \VUgras{касирка}
\textsuperscript{(14)}, що вносила \VUgras{рахунок} \textsuperscript{(17)} до
\VUgras{книги приїжджих} \textsuperscript{(16)}, попросила \VUgras{швейцара}
\textsuperscript{(18)} послати \VUgras{посильного} \textsuperscript{(19)} справити
моє доручення на його \VUgras{велосипеді} \textsuperscript{(20)}.

%%K t13-04-1 p
4. --- Я подякував їй і вийшов дати на чай \VUgras{носієві}
\textsuperscript{(24)} (чоловікові на всі руки), який привіз із вокзалу на своєму
\VUgras{ручному візку} \textsuperscript{(25)} мою \VUgras{капелюшну коробку}
\textsuperscript{(26)}, мою \VUgras{валізу} \textsuperscript{(27)} і мій
\VUgras{дорожній мішок} \textsuperscript{(28)}. \VUgras{Листоноша}
\textsuperscript{(21)}, що підійшов саме тієї хвилини, подав мені \VUgras{лист}
\textsuperscript{(22)}.

%%K t13-05-1 p
5. --- Я затримався ще трохи на порозі, коло \VUgras{поштової скриньки}
\textsuperscript{(23)}, дивлячись на рух по вулиці. \VUgras{Кучер}
\textsuperscript{(30)} \VUgras{омнібуса} \textsuperscript{(29)} вантажив на свого
воза \VUgras{скриню} \textsuperscript{(31)} одного \VUgras{пожильця}
\textsuperscript{(32)}, з яким прощався \VUgras{господар готелю}
\textsuperscript{(33)}. Молодий \VUgras{телеграфіст} \textsuperscript{(35)} ніс,
анітрохи не поспішаючи, \VUgras{депешу} \textsuperscript{(34)} одному з
пожильців; \VUgras{турист} \textsuperscript{(36)} з \VUgras{фотографічним
апаратом} \textsuperscript{(38)} через плече справлявся зі своїм
\VUgras{путівником} \textsuperscript{(37)}.

%%K t13-tit-1 sub
\VUsaut{4.0mm}
\VUpk{10.2pt}{40}{Година сніданку.}
\VUsaut{3.0mm}

%%K t13-06-1 p
6. --- Перед ґратами безжурний \VUgras{комісіонер} \textsuperscript{(39)} дрімав
між своїми \VUgras{гаками} \textsuperscript{(40)} і своєю \VUgras{скринькою
чистильника} \textsuperscript{(41)}, а молода \VUgras{пралька}
\textsuperscript{(42)} з \VUgras{кошиком} \textsuperscript{(43)} білизни сміялася
з поважного вигляду \VUgras{метрдотеля} \textsuperscript{(44)}, що стояв коло
дверей.

%%K t13-07-1 p
7. --- Вигляд \VUgras{кухаря} \textsuperscript{(45)}, який вийшов зі своєї
\VUgras{кухні} \textsuperscript{(47)} у \VUgras{підвалі} \textsuperscript{(48)},
щоб послати \VUgras{кухарчука} \textsuperscript{(46)} з дорученням, нагадав мені,
що близька година сніданку, і я пішов до ресторану на другому боці подвір'я, де
\VUgras{автомобіліст} \textsuperscript{(50)} ставив свій \VUgras{автомобіль}
\textsuperscript{(49)}, а \VUgras{механік} \textsuperscript{(52)} лагодив, за
вказівками \VUgras{велосипедиста} \textsuperscript{(53)}, \VUgras{бензиновий
трицикл} \textsuperscript{(51)}, який щойно зазнав аварії.

%%K t13-08-1 p
8. --- Я спитав у молодого \VUgras{грума} \textsuperscript{(54)}, який ніс
\VUgras{кухлі пива} \textsuperscript{(55)}, чи під верандою спільний стіл, і,
коли він сказав, що так, сів за один зі столиків, і мені подали чудовий
сніданок.

%%K t13-noto-2 noto
\VUnotes{143.2mm}{%
\textit{\VUgras{(*) За допомогою списку страв, поданого в словнику, учитель
легко може урізноманітнити наведене нижче меню.}}%
}

%%K t13-tit-2 sub
\VUsaut{4.0mm}
\VUpk{10.2pt}{40}{Добірний сніданок.}
\VUsaut{3.0mm}

%%K t13-09-1 p
9. --- Лакей приніс мені меню сніданку (*) і карту вин. Я вибрав і замовив на
\VUgras{закуску} \VUgras{креветок} з \VUgras{маслом}, а на першу страву ---
\VUgras{рубці по-канському}. \VUgras{Риби} я не брав, але попросив порцію
баранячої \VUgras{ноги} і, на \VUgras{городину}, \VUgras{шпинату} з вершками.
Потім, що жодна з \VUgras{солодких страв} мені не сподобалася, я велів подати
різні десерти: \VUgras{сир}, \VUgras{садовину} і \VUgras{печиво}. Я додав
чудового сент-емільйонського \VUgras{вина} і, дуже вдоволений своїм сніданком,
устав з-за столу, давши \VUgras{лакеєві} \textsuperscript{(57)} щедрі
\VUgras{чайові} \textsuperscript{(59)}.

%%K t13-10-1 p
10. --- Бачачи, що я збираюся йти, \VUgras{управитель} \textsuperscript{(60)}
запропонував мені вийти на балкон помилуватися чудовою панорамою
\VUgras{океану} \textsuperscript{(62)}. Удалині розрізнялися маленькі рибальські
\VUgras{човни} \textsuperscript{(63)}, \VUgras{вітрильники} \textsuperscript{(64)}
і величезний \VUgras{пакетбот} \textsuperscript{(65)}, чий чорний силует
вирізувався на білих \VUgras{хмарах} \textsuperscript{(67)} коло \VUgras{обрію}
\textsuperscript{(66)}. Легкий вітерець ворушив \VUgras{прапор}
\textsuperscript{(70)}, укріплений над \VUgras{вивіскою} \textsuperscript{(69)}
\VUgras{вестибюля} \textsuperscript{(68)}.

%%K t13-tit-3 sub
\VUsaut{3.0mm}
\VUpk{10.2pt}{40}{Розваги.}
\VUsaut{3.0mm}

%%K t13-11-1 p
11. --- Видовище було таке цікаве, що я вирішив другого дня піднятися на терасу
кав'ярні, взявши з собою \VUgras{бінокль} \textsuperscript{(71)} або
\VUgras{підзорну трубу} \textsuperscript{(72)}.

%%K t13-12-1 p
12. --- Пішовши з балкона, я подався до готельної кав'ярні випити
\VUgras{чогось} \textsuperscript{(73)} і зіграти партію на \VUgras{більярді}
\textsuperscript{(75)}. Я пройшов просто через залу кав'ярні, де багато
відвідувачів грало в \VUgras{шахи} \textsuperscript{(78)}, у \VUgras{шашки}
\textsuperscript{(79)}, у \VUgras{доміно} \textsuperscript{(80)}, у
\VUgras{триктрак} \textsuperscript{(81)} або в \VUgras{карти}
\textsuperscript{(82)}, і піднявся просто до \VUgras{більярдної}
\textsuperscript{(74)}.

%%K t13-13-1 p
13. --- Скінчивши партію, я спустився в нижній поверх і попросив у лакея
\VUgras{конверт} \textsuperscript{(84)}, \VUgras{паперу} \textsuperscript{(83)},
\VUgras{марку} \textsuperscript{(85)}, \VUgras{перо} \textsuperscript{(87)} і
\VUgras{чорнила} \textsuperscript{(86)}, щоб написати листа.

%%K t13-13-2 p
Потім я підкликав \VUgras{візника} \textsuperscript{(92)}, чий \VUgras{екіпаж}
\textsuperscript{(88)} спинився перед кав'ярнею. Я спитав його ціну за годину і
за кінець, і, коли ми сторгувалися, він відчинив мені \VUgras{дверцята}
\textsuperscript{(89)} і всадовив усередину. Він знову виліз на свої
\VUgras{козли} \textsuperscript{(91)}, жваво торкнув коней своїм \VUgras{батогом}
\textsuperscript{(93)}, і ми покотили на чудову прогулянку.
\VUsaut{6.0mm}
\VUfilet{22mm}
